
\documentclass[a4paper,12pt]{article}

\input{glyphtounicode}
\pdfgentounicode=1

\usepackage{mathptmx}
\usepackage{fancyhdr}
\usepackage{hyperref}
\usepackage{footmisc}
\usepackage[english]{babel}
\usepackage[utf8]{inputenc}
\usepackage[version=4]{mhchem}
\usepackage[usenames,dvipsnames]{color}
\usepackage[
  a4paper,
  top    = 2.00cm,
  bottom = 2.00cm,
  left   = 1.70cm,
  right  = 1.70cm
]{geometry}

\renewcommand{\baselinestretch}{1.25}
\renewcommand{\headrulewidth}{0pt}
\renewcommand{\footrulewidth}{0pt}
\renewcommand\footnotelayout{\fontsize{10}{10}\selectfont} 

\hypersetup{
    colorlinks=true,
    linkcolor=RoyalBlue,
    urlcolor=RoyalBlue,
    citecolor=RoyalBlue,
    anchorcolor=RoyalBlue
}

\urlstyle{same}
\pagestyle{fancy}
\fancyhf{}
% \fancyfoot[R]{\thepage}

\begin{document}

\begin{center}
  {\Huge \scshape {\fontsize{25}{30}\selectfont{Reza}} {\fontsize{25}{30}\selectfont{\textbf{Aliasgari Renani}}}} \\ 
  {\fontsize{12}{12}\selectfont 25/12/2025} $|$
  {\fontsize{12}{12}\selectfont Thesis Abstracts} $|$
  {\fontsize{12}{12}\selectfont TUDelft} $|$
  {\fontsize{12}{12}\selectfont PhD Position (2955)} $|$
  {\fontsize{12}{12}\selectfont Integrated Circuit Design}
\end{center}

\vspace{-20pt}
\noindent\rule{\textwidth}{0.5pt}

\noindent\textbf{MSc thesis Abstract}
\vspace{5pt}

Programmable microelectronic systems were evaluated under plasma-based radiation environments simulating space flight conditions. Several FPGA-based image-signal-processing pipelines were implemented on a Xilinx Zynq-7000 (Zybo Z7) board using manual Verilog coding and Simulink-generated HDL. Functionality was verified with automated testbenches. The designs were synthesized and validated in Vivado with HDMI I/O and live camera data. The board was then exposed to electron-beam plasma (25--60\,keV, up to 100\,mA) in low-pressure oxygen (10$^{-6}$--50\,Torr) to induce surface charging and X-rays, while undergoing thermal cycling (218--393\,K). System output was monitored through the HDMI output signal to detect functional disruptions, with board placement and cable routing arranged to preserve signal integrity during exposure. Initial runs maintained normal operation. 

\vspace{5pt}

The thesis is in progress: next steps are longer runs, higher electron energies within facility limits, and more fundamental characterization techniques to identify the intermittent effects before finalizing conclusions.

\vspace{25pt}

\noindent\textbf{BSc thesis Abstract}
\vspace{5pt}

This work is focused on investigation of electron and gamma quanta irradiation as part of cosmic radiation on microelectronic devices. Silicon oxide based and aluminum gated MOS structures fabricated on n-type silicon are investigated after irradiation with low energy electrons in a scanning electron microscope. The thermally stimulated current (TSC) technique in the temperature range 80\,K to 320\,K revealed a number of the electron-beam-induced charge traps. With the help of the capacitance-voltage (C-V) method, the traps revealed by the TSC were identified by their location (within dielectric, semiconductor, or at the interface) and by their nature (trap for electrons or for holes). Density of interface states and additional parameters were calculated and compared using both experimental techniques TSC and C-V. Further investigations were done on similar MOS structures after gamma irradiation in the temperature range 80\,K to 400\,K. The gamma-irradiation-induced traps were also identified by their location and nature. The dependency of peak temperature positions and peak magnitudes (resulting from the emission of charge carriers from energy levels) on temperature rate during a TSC scan was also investigated for gamma-irradiated MOS structures.



\end{document}
